\scanpage{177}
\mspage{492}
% position: top six pairs
\diagram{\begin{array}{l}
\text{dynamisme-équilibre}\\
\text{élan-assise}\\
\text{ardeur-persévérance}\\
\text{fougue- patience}\\
\text{passion-sérénité}\\
\text{tenacité-détachement.}
\end{array}}
% position: upper paragraph
J'y joindrais bien encore les deux couples suivants, parmi une dizaine de
``retardataires'' qui me sont venus encore ce matin, sur la lancée de ma réflé-
xion de hier:
% position: upper two pairs
\diagram{\begin{array}{l}
\text{savoir-connaître}\\
\text{expliquer-comprendre.}
\end{array}}
% position: middle paragraph
Est-il besoin de préciser que dans ces couples, c'est le terme ``yang''
ou ``masculin'' qui est mis le premier, suivant l'usage de notre société patriar-
cale, où l'homme donne le nom au couple? Par contre, alors que la société chi-
noise traditionnelle est considérablement plus patriarcale encore que la nôtre,
quand on suit l'usage chinois pour parler de la relation du yin et du yang,
on met toujours en premier le yin (``féminin''), en parlant p. ex. d'``équilibre
yin-yang'' (au lieu de yang-yin). Le sens de cet usage est sûrement dans l'intui-
tion-archétype que c'est le yang qui naît du yin, lequel est le principe ``le
plus primitif'' des deux, et non l'inverse...

% position: lower middle paragraph
Ce n'est pas le lieu ici de me lancer dans des comentaires sur l'un ou
l'autre de ces couples. Pour le lecteur qui ne ``sentirait rien'' en les voyant,
ce serait de toutes façons peine perdue; et celui qui se sent interpellé par
eux, qui sent (fut-ce obscurément) que chacun d'eux a quelque chose à lui di-
re sur le monde et sur lui-même - sur l'équilibre et le déséquilibre, sur la
dynamique interne des êtres et des choses..., celui-là peut se passer de com-
mentaires circonstanciés, et prendre cette interpellation comme un point de
départ pour sa propre réflexion.

% position: lower handwritten number and paragraph
\margin{$(111^1)$}
Il y a un point seulement sur lequel je voudrais insister ici, commun à
tous les ``couples'' yin-yang sans exception. C'est la chose aussi la plus cru-
ciale de toutes, il me semble, pour une compréhension de la nature de la rela-
tion entre le yin et le yang, et par là, de la nature de chacun de ces deux
principes (ou énergies, ou aspects, ou forces...) dans l'Univers. C'est ceci:
chacun des deux termes d'un de ces couples, tel action-inaction, \mbox{e n\quad l' a b-}
% position: lower left short stroke
\margin{\unsure{-}}
% position: foot pencil
\margin{176}
